\chapter{代数式}

本章覆盖知识点如下：
\begin{introduction}
  \item 代数式
  \item Ask for help
  \item Optimization Problem
  \item Property of Cauchy Series
  \item Angle of Corner
\end{introduction}


\section{整体代入思维}


\begin{definition}[可积性] \label{def:int} 
设 $ f(x)=\sum\limits_{i=1}^{k} a_i \chi_{A_i}(x)$ 是 $E$ 上的\textbf{非负简单函数}，中文其中 $\{A_1,A_2,\ldots,A_k\}$ 是 $E$ 上的一个可测分割，$a_1,a_2,\ldots,a_k$ 是非负实数。定义 $f$ 在 $E$ 上的积分为 $\int_{a}^b f(x)$
\begin{equation}
   \label{inter}
   \int_{E} f dx = \sum_{i=1}^k a_i m(A_i) \pi \alpha\beta\sigma\gamma\nu\xi\epsilon\varepsilon. \oint_{a}^b\ointop_{a}^b\prod_{i=1}^n
\end{equation}
一般情况下 $0 \leq \int_{E} f dx \leq \infty$。若 $\int_{E} f dx < \infty$，则称 $f$ 在 $E$ 上可积。
\end{definition}

\subsection{整体代入-题目1}

已知三个数 \(a, b, c\)，用 \(\max\{a, b, c\}\) 表示这三个数中最大的数，如果 \(\dfrac{1}{3}(a+b+c)=\max\{a,b,c\}\)，那么：

(1) 当 \(a=3,\ b=2+x,\ c=2x+1\) 时，\(x\) 的值是 \_\_\_\_\_\_；

(2) 当 \(a=\dfrac{x+y}{2},\ b=\dfrac{y+z}{3},\ c=\dfrac{z+x}{7}\)，且 \(x+2y+3z=26\) 时，\(x+y+z\) 的值是 \_\_\_\_\_\_。

\section{代数最值问题}

\begin{proposition}[多知识点融合] \label{pro:max}
【知识点】代入消元思想、线性函数的值域、线性规划

【分析】XXX.
\end{proposition}

\subsection{代数最值问题-题目1}

\(A, B, C\) 三种原料每袋的重量（单位：kg）依次是 1, 2, 3，每袋的价格（单位：万元）依次是 3, 2, 5。现生产某种产品需要 \(A, B, C\) 这三种原料的袋数依次为 \(x_1, x_2, x_3\)（\(x_1, x_2, x_3\) 均为正整数），则生产这种产品时需要的这三类原料的总重量 \(W\)（单位：kg）= \_\_\_\_\_\_（用含 \(x_1, x_2, x_3\) 的代数式表示）；为了提升产品的品质，要求 \(W \geq 13\)，当 \(x_1, x_2, x_3\) 的值依次是 \_\_\_\_\_\_ 时，这种产品的成本最低。

（答案见第~\pageref{代数最值问题-题目1}~页）


\subsection{代数最值问题-题目2}



【阅读材料】：

材料一：对于实数 \(x, y\) 定义一种新运算 \(K\)，规定：\(K(x, y) = ax + by\)（其中 \(a, b\) 均为非零常数），等式右边是通常的四则运算。比如：\(K(1, 2) = a + 2b\)；\(K(-2, 3) = -2a + 3b\)。

已知：\(K(1, 2) = 7\)；\(K(-2, 3) = 0\)。

材料二：“已知 \(x, y\) 均为非负数，且满足 \(x + y = 8\)，求 \(2x + 3y\) 的范围”，有如下解法：

\[\because x + y = 8, \quad \therefore x = 8 - y,\]
\[\because x, y \text{ 是非负数，} \therefore x \geq 0 \text{ 即 } 8 - y \geq 0, \quad \therefore 0 \leq y \leq 8,\]
\[\because 2x + 3y = 2(8 - y) + 3y = 16 + y, \quad \therefore 16 \leq 16 + y \leq 24, \quad \therefore 16 \leq 2x + 3y \leq 24.\]

【回答问题】：

(1) 求出 \(a, b\) 的值；

(2) 已知 \(x, y\) 均为非负数，\(x + 2y = 10\)，求 \(4x - y\) 的取值范围；

(3) 已知 \(x, y, z\) 都为非负数，\(K(y, z) = 3 + x, \quad K\left(x, \frac{1}{2}y\right) = 4 - 3x\)，求 \(x - 3y + 4z\) 的取值范围。


（答案见第~\pageref{代数最值问题-题目2}~页）



\subsection{代数最值问题-题目3-2023七下·石景山期末}

对于二元一次方程的任意一个解给出如下定义：若 \(x \geq y\)，则称 \(x - y\) 为方程的“关联值”；若 \(x < y\)，则称 \(y - x\) 为方程的“关联值”。已知方程 \(x + y = 3\)。

\begin{enumerate}
    \item 写出方程 \(x + y = 3\) 的一个解，并指明此时方程的“关联值”；
    \item 若“关联值”为 \(4\)，写出所有满足条件的解；
    \item 直接写出方程 \(x + y = 3\) 的最小“关联值”；当“关联值”为 \(1\) 时，直接写出 \(x\) 的取值范围。
\end{enumerate}

（答案见第~\pageref{代数最值问题-题目3-2023七下·石景山期末}~页）


\subsection{代数最值问题-题目3-2025北京三帆中学初一下期中第2题}

已知关于 \(x, y\) 的方程组
\[
\begin{cases}
x + 2y = m + 5 \\
2x + y = m + 4
\end{cases}
\]
其中 \(x, y\) 为非负数，若 \(S = x - y\)，则 \(S\) 的取值范围是 \_\_\_\_\_\_\_\_\_\_。


（答案见第~\pageref{代数最值问题-题目3-2025北京三帆中学初一下期中第2题}~页）
